%!TEX root = ../template.tex
%%%%%%%%%%%%%%%%%%%%%%%%%%%%%%%%%%%%%%%%%%%%%%%%%%%%%%%%%%%%%%%%%%%%
%% abstrac-pt.tex
%% ISEL thesis document file
%%
%% Abstract in Portuguese
%%%%%%%%%%%%%%%%%%%%%%%%%%%%%%%%%%%%%%%%%%%%%%%%%%%%%%%%%%%%%%%%%%%%
A presente dissertação enquadra-se na área da Engenharia Informática e da Saúde, com foco no desenvolvimento de sistemas de recomendação baseados em Grafos de Conhecimento para potenciar descobertas de novos fármacos e investigação clínica. O elevado volume de dados biomédicos atualmente disponível é frequentemente não uniforme, fragmentado e não estruturado, o que apresenta um desafio para a sua integração e análise de forma eficaz. O objetivo principal é desenvolver um sistema capaz de prever novas hipóteses biomédicas através da predição de ligações em grafos multi-relacionais, com o objetivo de identificar oportunidades de reposicionamento de fármacos. A investigação procura, assim, analisar de que forma a diversidade na tipologia de entidades e relações representadas no grafo, explorada através de modelos como as Redes Neuronais Convolucionais Relacionais (R-GCN), contribui para o aumento da qualidade das recomendações, quando comparada com abordagens baseadas em grafos onde não há distinção dos tipos de entidades nem de relações.

A metodologia proposta prevê a construção de um Grafo de Conhecimento através da integração de dados regulatórios, utilizando a framework MedJsonify como ferramenta de apoio à extração de entidades e ao seu mapeamento para ontologias de referência. A eficácia da solução será avaliada quantitativamente recorrendo às métricas Mean Reciprocal Rank (MRR) e Hits@k, permitindo avaliar a precisão do modelo na predição de novas ligações. Adicionalmente, será realizada uma validação qualitativa através da comparação das hipóteses geradas com evidências presentes na literatura científica e registos de ensaios clínicos. Espera-se que a aprendizagem sobre estas estruturas relacionais permita identificar padrões clínicos relevantes, apoiando a descoberta de novas indicações terapêuticas.

\begin{keywords}
Grafos de Conhecimento, Redes Neuronais Convolucionais Relacionais (R-GCN), Predição de Ligações, Reposicionamento de Fármacos.
\end{keywords}
